% Source: https://artofproblemsolving.com/community/c3h2193036_math_problem
% Stable identity: aops-topic-2193036-q01
Number the minutes $1,2,3,\dots$. The guard on duty alternates, so the guards on duty in minutes $t$ and $t+k$ are the same if $k$ is even and different if $k$ is odd.

\textbf{(a)} The robber is at $3$ at the end of minute $3$ and at $5$ at the end of minute $5$. Minutes $3$ and $5$ have the same guard on duty. If it is the guard at $3$, the robber is caught at $3$; if it is the guard at $5$, he is caught at $5$. So the probability is $\boxed{0}$. (Multiplying $\dfrac{1}{2}\cdot\dfrac{1}{2}=\dfrac{1}{4}$ would be wrong: the two events ``no guard at $3$ in minute $3$'' and ``no guard at $5$ in minute $5$'' are not independent; they never happen together.)

\textbf{(b)} The robber is never in danger at $0$, $1$, $2$ or $4$.
\begin{itemize}
\item \emph{At $3$.} Let $T$ be the minute in which he first reaches $3$. He survives only if the guard at $5$ is on duty in minute $T$; this has probability $\dfrac{1}{2}$, because the guards' starting coin is independent of the robber's coin flips. Then the guard at $3$ is on duty in minute $T+1$, so he must move on to $4$ in minute $T+1$: probability $\dfrac{1}{2}$.
\item \emph{At $4$.} Suppose he is at $4$ at the end of exactly $m\ge1$ minutes, $T+1,\dots,T+m$, and reaches $5$ in minute $T+m+1$; this has probability $\left(\dfrac{1}{2}\right)^{m-1}\cdot\dfrac{1}{2}=\left(\dfrac{1}{2}\right)^m$. The guard at $5$ was on duty in minute $T$, so in minute $T+m+1$ the guard at $3$ is on duty exactly when $m+1$ is odd, that is, when $m$ is even.
\item \emph{At $5$.} He must then move on to $6$ in the very next minute (when the guard at $5$ is on duty): probability $\dfrac{1}{2}$.
\end{itemize}
Adding over the even values $m=2,4,6,\dots$,
\[P=\dfrac{1}{2}\cdot\dfrac{1}{2}\cdot\left(\dfrac{1}{4}+\dfrac{1}{16}+\dfrac{1}{64}+\cdots\right)\cdot\dfrac{1}{2}=\dfrac{1}{8}\cdot\dfrac{\dfrac{1}{4}}{1-\dfrac{1}{4}}=\dfrac{1}{8}\cdot\dfrac{1}{3}=\boxed{\dfrac{1}{24}}.\]
